\section{统一设置、覆盖顺序与作用域}
\subsection{统一设置与图形入口}
\key{\gradbarssetup{参数列表}} 设置当前作用域的默认值；
\key{\gradbar[参数列表]{数值}} 绘制一个数据条，其中可选参数只覆盖本次绘制。
数值与最大值可用十进制或 e/E 科学计数数字，不接受逗号、单位或计算表达式。

\begin{example}{示例 1：统一设置与单次覆盖}
\gradbarssetup{width=23mm,
  max=100,precision=1,theme=blue}
\gradbar{72.4}

\gradbar[theme=teal]{92.7}

\gradbar{81.6}
\end{example}

\subsection{作用域}
设置遵循普通 TeX 分组规则。在导言区调用会影响整篇文档；
在一对花括号中调用只影响该组。表格单元格也有自己的分组，
因此整列共用的设置应放在表格环境之前。

\begin{example}{示例 2：局部设置不会泄漏}
{\gradbarssetup{theme=gray,max=10}
 \gradbar{8}}

\gradbar{80}
\end{example}
